
\section{Related Work}

\subsection{Low-Rank Matrix Factorization}

\textbf{Classical SVD.} The Eckart--Young theorem (Eckart \& Young, 1936) states 
that the best rank-$r$ approximation of $W$ in Frobenius norm is the truncated 
SVD. MinimalSVD builds on this foundation.

\textbf{Randomized SVD.} Halko et al. (2011) provide randomized algorithms with 
complexity $\mathcal{O}(dr \log k)$. For our layer sizes, exact SVD is 
feasible within seconds.

\textbf{Weighted low-rank.} Golub et al. (1987) and Srebro \& Jaakkola (2003) study 
weighted variants. Activation-aware methods (Lin et al., 2023; Xiao et al., 
2023) weight the factorization by activation statistics. MinimalSVD uses 
\textbf{unweighted} SVD but \textbf{adaptive rank}, offering a complementary axis of 
improvement.

\textbf{Tensor decomposition.} Novikov et al. (2015) use tensor-train 
decomposition. MinimalSVD uses matrix SVD, a simpler primitive.

\subsection{Rank Selection Strategies}

\textbf{Fixed rank.} Most low-rank methods use a fixed global rank (Yu et al., 
2017). Simple but ignores heterogeneity.

\textbf{Adaptive rank in LoRA.} AdaLoRA (Zhang et al., 2023) learns per-layer 
rank budgets. MinimalSVD uses \textbf{spectral energy}, a fixed criterion 
without learning.

\textbf{Energy-based truncation.} Standard in signal processing. MinimalSVD 
applies it to transformer attention projections.

\subsection{Parameter-Efficient Fine-Tuning}

\textbf{LoRA.} Hu et al. (2021) freeze base weights and train rank-$r$ adapters. 
MinimalSVD uses LoRA as the recovery mechanism after compression.

\textbf{DoRA.} Liu et al. (2024) decompose weights into magnitude and direction. 
Orthogonal to our work.

\subsection{Attention Variants}

\textbf{Standard MHSA.} Vaswani et al. (2017) introduce the four-projection 
attention. MinimalSVD preserves the four-projection structure but factorizes 
each.

\textbf{GQA.} Ainslie et al. (2023) share KV heads. MinimalSVD supports GQA via 
\texttt{repeat\_kv}.

\textbf{MQA.} Shazeer (2019) uses a single KV head. Compatible with our approach.

\subsection{Model Compression}

\begin{table}[H]
\centering
\small
\begin{tabular}{lllll}
\toprule
\textbf{Method} & \textbf{Type} & \textbf{Ratio} & \textbf{Retrain} & \textbf{Calibration} \\
\midrule
DistilBERT & Distillation & 1.6$$\times$ & Yes & -- \\
TinyBERT & Distillation & 7$$\times$ & Yes & -- \\
SparseGPT & Pruning & 2-4$$\times$ & No & Required \\
Wanda & Pruning & 2$$\times$ & No & Required \\
AWQ & Quantization & 4$$\times$ & No & Required \\
GPTQ & Quantization & 4$$\times$ & No & Required \\
A3 & Low-rank + diagonal Fisher & 5$$\times$ & LoRA & Required \\
AASVD & Low-rank + full covariance & 5$$\times$ & LoRA & Required \\
\textbf{MinimalSVD} & \textbf{Low-rank + energy} & \textbf{50\%$$\times$} & \textbf{LoRA} & \textbf{Not required} \\
\bottomrule
\end{tabular}
\end{table}

\subsection{Position of Our Work}

MinimalSVD differs from prior work along three axes:

\begin{table}[H]
\centering
\small
\begin{tabular}{lll}
\toprule
\textbf{Axis} & \textbf{Prior} & \textbf{MinimalSVD} \\
\midrule
\textbf{Rank selection} & Fixed or learned & Spectral energy \\
\textbf{Calibration} & Often required & Not required \\
\textbf{Code complexity} & High (framework) & Minimal (auditable) \\
\bottomrule
\end{tabular}
\end{table}

\textbf{Key distinction.} MinimalSVD is designed as a \textbf{reference baseline}. It 
uses the same energy-adaptive rank criterion as MinimalNOVA and MinimalPRISM, 
but without the injection-time safety protocols and whole-model traversal. 
This makes the code shorter and easier to audit, at the cost of requiring 
the caller to handle model traversal.



\subsection{Comparison of the Minimal Family}

\begin{lstlisting}

---
\end{lstlisting}